\documentclass[11pt]{article}
\usepackage[margin=1in]{geometry}
\usepackage{amsmath,amssymb}

\title{Natural-Metric Newton Option (Short Note)}
\author{}
\date{}

\begin{document}
\maketitle

\section*{Overview}
For variational inference with Gaussian parameters
\[
\theta = \begin{bmatrix}\mu \\ \ell\end{bmatrix},
\quad \ell = \log \sigma,
\]
we added a Newton option with \texttt{newton\_metric = "natural"}.
This keeps the same objective and Hessian estimate, but applies Newton in
approximately Fisher-whitened coordinates.

\section*{Metric Scaling Used}
For a diagonal Gaussian family, the implementation uses a block-diagonal scale
\[
S = \mathrm{diag}\!\left(\sigma_1,\dots,\sigma_d,\sqrt{\tfrac12},\dots,\sqrt{\tfrac12}\right).
\]
This is consistent with the diagonal Fisher factors (up to constants) for
\(\mu\) and \(\ell\) blocks.

\section*{Where This Scaling Comes From}
For one scalar Gaussian factor \(q(z)=\mathcal N(\mu,\sigma^2)\), with
\(\ell=\log \sigma\),
\[
\log q(z) = -\ell - \frac12\left(\frac{(z-\mu)^2}{\sigma^2} + \log(2\pi)\right).
\]
The score components are
\[
\partial_\mu \log q = \frac{z-\mu}{\sigma^2},
\qquad
\partial_\ell \log q = -1 + \frac{(z-\mu)^2}{\sigma^2}.
\]
Using \(z\sim q\), the Fisher matrix entries are
\[
F_{\mu\mu} = \mathbb E\!\left[(\partial_\mu \log q)^2\right] = \frac{1}{\sigma^2},
\quad
F_{\ell\ell} = \mathbb E\!\left[(\partial_\ell \log q)^2\right] = 2,
\quad
F_{\mu\ell}=0.
\]
For independent dimensions, \(F\) is block diagonal:
\[
F = \mathrm{diag}\!\left(\frac{1}{\sigma_1^2},\dots,\frac{1}{\sigma_d^2},2,\dots,2\right),
\]
so
\[
F^{-1/2}
= \mathrm{diag}\!\left(\sigma_1,\dots,\sigma_d,\sqrt{\tfrac12},\dots,\sqrt{\tfrac12}\right)
= S.
\]
That is exactly the scaling used in the implementation.

\section*{Newton Step in Natural-Metric Coordinates}
Given gradient \(g\) and Hessian \(H\) in the original coordinates, we form
\[
\tilde g = S g,\qquad \tilde H = S H S.
\]
Then solve the regularized Newton system
\[
\tilde H\,\tilde\Delta = \tilde g,
\]
and map back
\[
\Delta = S \tilde\Delta.
\]
The code applies SPD projection/regularization to \(\tilde H\) before solving.

\section*{Why This Newton Step Is Reasonable}
The local KL geometry is determined by \(F\). A second-order local model in
original coordinates is
\[
m(\Delta\theta) = g^\top \Delta\theta + \frac12 \Delta\theta^\top H \Delta\theta.
\]
With the change of variables \(\Delta\theta = S u\) (Fisher whitening),
\[
m(u) = (Sg)^\top u + \frac12 u^\top (S H S) u,
\]
which is a standard Euclidean quadratic model in \(u\). Newton in \(u\)-space
therefore gives the transformed system above and maps back with
\(\Delta\theta = S\tilde\Delta\).

This is an approximate natural-Newton method (practical and common): it uses
the Fisher metric as a preconditioning geometry but does not include full
Riemannian connection terms (\(\partial F/\partial\theta\)).

\section*{Trust Region Variant}
For trust-region Newton, the same transform is used:
\[
(\tilde g,\tilde H) = (Sg,SHS),
\]
the trust-region subproblem is solved in transformed coordinates, and the step is
mapped back with \(\Delta = S\tilde\Delta\). Radius adaptation and acceptance
logic remain unchanged.

\section*{Configuration}
Use:
\begin{verbatim}
optimizer_method = "newton"
optimizer_config = VINewtonOptimizerConfig(
    newton_metric="natural",
    ...
)
\end{verbatim}
The default remains \texttt{newton\_metric = "standard"}.

\section*{Note}
This is an approximate natural-Newton construction (a practical preconditioned
Newton), not a full Riemannian Newton method with connection terms.

\end{document}
